\documentclass[11pt, a4paper]{article}

\usepackage[utf8]{inputenc}
\usepackage{geometry}
\geometry{a4paper, margin=1in}
\usepackage{hyperref}
\usepackage{graphicx}
\usepackage{cite}

\title{Decentralized Point of Interest Prediction: A Privacy-Preserving Edge Architecture}
\author{Dein Name \and Name deines Partners}
\date{\today}

\begin{document}

\maketitle

\begin{abstract}
\begin{itemize}
    \item Das Problem (Privacy bei Location-based Services)
    \item Der Lösungsansatz (Decentralized/Edge-based Prediction, K-Means User Clustering)
    \item Kurzer Ausblick auf die Ergebnisse
\end{itemize}
\end{abstract}

\section{Introduction}
\begin{itemize}
    \item Motivation: Warum brauchen wir Recommender Systems für Points of Interest (POIs)?
    \item Problem: Klassische Ansätze sammeln komplette Location Histories zentral (Privacy Issue).
    \item Proposed Solution: Eine dezentrale Architektur, die auf dem Endgerät inferiert.
    \item Paper Structure: Kurzer Überblick, was in den nächsten Sektionen passiert.
\end{itemize}

\section{Related Work}
\subsection{Point of Interest Recommendation Systems}
\begin{itemize}
    \item Historischer Überblick über POI Recommender (Collaborative Filtering, Matrix Factorization).
    \item State of the Art für Time-based und Location-based Empfehlungen.
\end{itemize}

\subsection{Federated Learning and Decentralized Approaches}
\begin{itemize}
    \item Überblick über Federated Learning.
    \item Wie unterscheidet sich unser Ansatz (Centroids on device) von klassischem Federated Learning (wo Gewichte hin- und hergeschickt werden)?
    \item Vorteile der Edge-Inferenz in Bezug auf Privacy.
\end{itemize}

\section{Methodology}
\subsection{Dataset and Data Preprocessing}
\begin{itemize}
    \item Beschreibung des Foursquare Datensatzes (TSMC2014).
    \item Data Cleaning: Herausfiltern von irrelevanten Orten (Airports, Train Stations, Home) für ein "Tourist Scenario".
    \item Mapping der rohen Foursquare-Kategorien auf Level-1 Oberkategorien (z.B. "Dining and Drinking").
\end{itemize}

\subsection{Decentralized Architecture Overview}
\begin{itemize}
    \item High-Level Architektur (Big Picture).
    \item Was passiert in der Cloud (Training von Baseline, Clustering auf Public Users)?
    \item Was passiert on-device (Cluster-Matching, Inferenz)?
\end{itemize}

\subsection{User Partitioning via K-Means Clustering}
\begin{itemize}
    \item Konzept: Clustering von "Public Users" in der Cloud basierend auf ihrer Kategorie-Verteilung.
    \item Anstatt Location History zu speichern, extrahieren wir nur die Interessen (Prozentuale Verteilung der Level-1 Kategorien).
    \item K-Means trainiert k Centroids.
    \item Device-Side: Das Endgerät lädt nur die k Centroids herunter, ordnet das lokale (private) Profil einem Cluster zu, ohne die eigenen Daten in die Cloud zu schicken.
\end{itemize}

\subsection{Time-Based Category Prediction}
\begin{itemize}
    \item Das Baseline-Modell: Eine einfache, schnelle Lookup-Tabelle ($\mathcal{O}(1)$).
    \item Training erfolgt extrem datensparsam: Wir speichern keinen User-Context und keine Geodaten, sondern \textit{nur} die (Cluster, Hour) $\rightarrow$ Top-Kategorien Map.
    \item Das Modell wird komplett auf dem Endgerät (Edge) deployed.
    \item Das Endgerät nutzt die eigene lokale Uhrzeit und den eigenen (lokal berechneten) Cluster, um die nächste "Kategorie" vorherzusagen.
\end{itemize}

\subsection{Advanced Prediction Model (Dedicated Approach)}
\begin{itemize}
    \item Hier kommt der Part deines Partners hin!
    \item Komplexeres Modell (z.B. Deep Learning, Graph Neural Networks, etc.).
    \item Wie integriert sich dieses in die dezentrale Architektur?
\end{itemize}

\section{System Implementation and Live Retrieval}
\begin{itemize}
    \item Sobald das Modell on-device z.B. "American Restaurant" vorhergesagt hat, muss ein echter Ort her.
    \item Implementation: Edge-Request an die Foursquare Places API (Proxy-Backend).
    \item Query-Strategien (Popularity vs. Distance vs. Relevance).
\end{itemize}

\section{Limitations}
\subsection{Dataset Constraints}
\begin{itemize}
    \item Geografischer Bias: Unser Modell ist aktuell stark auf die Charakteristiken von Großstädten (New York City und Tokyo) begrenzt.
    \item Dichte: NYC hat extrem viele Foursquare-Daten, in ländlichen Regionen würde die Time-based Prediction anders aussehen.
\end{itemize}

\subsection{Privacy Limitations of the API Retrieval}
\begin{itemize}
    \item Der Trade-Off: Obwohl wir die Vorhersage (welche Kategorie will der User?) lokal auf der Edge machen und keine Location History zentral speichern, müssen wir für den \textit{Retrieval} (Welche Orte gibt es hier?) einen API-Call machen.
    \item Das erfordert das Senden der aktuellen Koordinaten an Foursquare. 
    \item Möglicher Ausblick: Lokale Offline-Kartenbanken (OSM) auf dem Handy, um auch den Retrieval komplett offline und 100\% private zu machen.
\end{itemize}

\section{Evaluation and Results}
\begin{itemize}
    \item Evaluierungsmetriken (Hit Rate @ K, Accuracy @ 1).
    \item Splitting-Strategie (Chronologischer Split, um Data Leakage in die Zukunft zu verhindern).
    \item Vergleich: Globales Time-based Model vs. Unser Cluster-spezifisches Time-based Model.
    \item (Hier fließen dann auch die Ergebnisse vom Modell deines Partners ein).
\end{itemize}

\section{Conclusion and Future Work}
\begin{itemize}
    \item Was haben wir erreicht? (Ein funktionierendes, Privacy-freundliches Edge-Recommender System).
    \item Fazit: Ist Dezentralisierung für Location-based Services machbar?
    \item Ausblick auf zukünftige Forschung (z.B. Integration von Offline-POI-Datenbanken).
\end{itemize}

\bibliographystyle{ieeetr}
\begin{thebibliography}{99}

\bibitem{foursquare_dataset}
Dingqi Yang et al., 
\textit{Modeling Spatial-Temporal Contexts for POI Recommendation},
IEEE Transactions on Knowledge and Data Engineering, 2014.

% Füge hier deine weiteren Paper zu FL und POI Recommendation ein

\end{thebibliography}

\end{document}
